%% kapa: exam=1 round=28 year=2017 period=2 form=A subject=appraisal_law subject_ko=감정평가관계법규 seq=21 number=21
%% source: Q-Net 감정평가사 기출문제 (cst00309) artlSeq=5199603 file=extracted/1008/5199603/2017년 제28회 감정평가사 1차시험 문제지/2교시 A형.pdf
%% conversion: kordoc 4.24.51 -> tools/gen1.py (auto); review=proofread-sonnet; flags=none
\begin{question}{21}
국유재산법령상 일반재산의 대부에 관한 설명으로 옳은 것은?
\begin{choices}
  \item 일반재산은 대부는 할 수 있으나 처분은 할 수 없다.
  \item 영구시설물의 축조를 목적으로 하는 토지와 그 정착물의 대부기간은 50년이 넘도록 정할 수 있다.
  \item 대부기간이 끝난 일반재산에 대하여 종전의 대부계약을 갱신할 수 있는 경우에도 수의계약의 방법으로 대부할 수 있는 경우에는 1회만 갱신할 수 있다.
  \item 중앙관서의 장등은 연간 대부료의 일부를 대부보증금으로 환산하여 받아야 한다.
  \item 일반재산을 주거용으로 대부계약을 하는 경우에는 수의(隨意)의 방법으로 대부계약의 상대방을 결정할 수 있다.
\end{choices}
\end{question}
\answer{5}
